\section{Cycle operations and the two inputs}\label{sec:cycles}

A correspondence $Z\in\operatorname{CH}^a(P\times Q)_{\mathbb Q}$ between smooth projective varieties acts by
\begin{equation}\label{eq:degree-convention}
 Z_*x=(p_Q)_*\bigl(\operatorname{cl}(Z)\cup p_P^*x\bigr):
 H^j(P)\longrightarrow H^{j+2e}(Q)(e),\qquad e=a-\dim P.
\end{equation}
The twist makes this a map of Hodge structures of weight $j$. Composition through a variety of dimension $r$ adds codimensions and subtracts $r$; the cycle class of the composition acts as the composite \cite{Fulton}. In particular a degree-preserving correspondence from a surface to degree two has codimension two.

\begin{lemma}[Weight-one operations]\label{cyc:h1}
For complex abelian varieties $P,Q$, pullback identifies
\[
 \operatorname{Hom}(Q,P)\otimes\mathbb Q
 \simeq
 \operatorname{Hom}_{\mathrm{HS}}\bigl(H^1(P,\mathbb Q),H^1(Q,\mathbb Q)\bigr).
\]
Thus arbitrary rational weight-one Hodge maps can act algebraically and independently on the two factors of an ordered tensor product of first cohomology groups.
\end{lemma}

\begin{proof}
Dualize a rational Hodge map to a map on first homology. After multiplication by a positive integer it preserves the integral lattices. The Hodge condition says that its real extension commutes with the complex structures on the universal covers. It therefore induces a holomorphic group homomorphism $Q\to P$, which is algebraic by projectivity. Conversely pullback by every such homomorphism is a Hodge map. This is the usual equivalence between abelian varieties up to isogeny and polarizable rational Hodge structures of weight one, written contravariantly \cite[Theorem 6.8 and Corollary 6.9]{Milne}. Transpose graphs realize the pullbacks. Products of two such correspondences, and their rational linear combinations, act independently on the two tensor factors. Preservation of the chosen polarizations is not required.
\end{proof}

\begin{lemma}[Degree projectors]\label{cyc:projectors}
Every cohomological K\"unneth projector on a product of abelian varieties is induced by a rational algebraic correspondence. If $U=H^1(B,\mathbb Q)$ and $m:B\times B\to B$ is addition, then the mixed part of $m^*$ gives an injection
\[
 H^2(B,\mathbb Q)=\bigwedge\nolimits^2U\longrightarrow U\otimes U,
 \qquad u\wedge v\longmapsto u\otimes v-v\otimes u.
\]
\end{lemma}

\begin{proof}
For an abelian variety $B$ of dimension $b$, multiplication by $2$ acts on $H^j(B)$ by $2^j$. Thus the rational polynomial
\[
 e_j(T)=\prod_{\substack{0\le k\le 2b\\k\ne j}}
 \frac{T-2^k}{2^j-2^k}
\]
applied to the correspondence $[2]^*$ is the projector onto $H^j(B)$. Apply these polynomials separately to each factor of a product. Finally $m^*u=p_1^*u+p_2^*u$ on $H^1$; multiplication in the exterior algebra, with its usual signs, gives the displayed formula.
\end{proof}

\begin{lemma}[Rational descent of a cycle span]\label{cyc:descent}
Let $E$ be a finite-dimensional rational vector space of cohomological maps, and let $E_{\mathrm{alg}}\subset E$ be the image of rational algebraic correspondences of a fixed degree. For every field extension $K/\mathbb Q$,
\[
 E\cap(E_{\mathrm{alg}}\otimes_{\mathbb Q}K)=E_{\mathrm{alg}}
 \quad\text{inside }E\otimes_{\mathbb Q}K.
\]
The same assertion holds if only a specified rational subspace of the source is considered.
\end{lemma}

\begin{proof}
The cycle class map and restriction of its action are rational linear maps, so $E_{\mathrm{alg}}$ is a rational subspace. If a rational target lies in the scalar-extended span of finitely many rational cycle actions, write these actions as columns in a rational basis of $E$. The linear system with that target as right-hand side is consistent over $K$, and row reduction over $\mathbb Q$ gives a rational solution. Only finitely many cycles are used.
\end{proof}

\begin{proposition}[Correspondences and surface covers]\label{cyc:equivalence}
Conditions \textup{(i)} and \textup{(ii)} of Theorem~\ref{main:generic} are equivalent. The implication from the surface condition to a nonzero correspondence on $V$ does not require the full Hodge-group hypothesis.
\end{proposition}

\begin{proof}
Suppose first that $f:Y\to S$ satisfies \eqref{eq:wedge-condition}. Let $D=\operatorname{Alb}(Y)$, let $a:Y\to D$ be an Albanese map, and put $b=\dim D$. Albanese pullback identifies holomorphic one-forms \cite[Chapter VI, \S9.2, pp.~315--316]{Albanese}. Hence there is
\[
 \beta\in H^{2,0}(D)=\bigwedge\nolimits^2H^{1,0}(D),\qquad
 y=f_*a^*\beta\ne0.
\]
In particular $b\ge2$ and $y\in H^{2,0}(S)\subset V_{\mathbb C}$. The surface correspondence $(a,f)_*[Y]$ has codimension $b$ on $D\times S$. Its transpose acts by
\[
 a_*f^*:H^2(S)\longrightarrow H^{2b-2}(D)(b-2).
\]
This action is nonzero on $V$: the projection formula gives
\begin{equation}\label{eq:albanese-nonzero}
 \int_D\beta\cup a_*f^*\overline y
 =\int_Ya^*\beta\cup f^*\overline y
 =\int_Sy\cup\overline y\ne0.
\end{equation}

We give an algebraic degree adjustment. For an ample line bundle $L$ on $D$, written additively, set
\[
 E=m^*L-p_1^*L-p_2^*L\in\operatorname{CH}^1(D\times D),
\]
where $m$ is addition. The class of $E$ is purely mixed. Put $U=H^1(D,\mathbb Q)$ and write this class as $\sum_i u_i\otimes v_i$, choosing a basis $(u_i)$ of $U$. Its matrix
\[
 M:U^\vee\longrightarrow U,\qquad
 M(\varphi)=\sum_i\varphi(u_i)v_i
\]
is invertible: it is the alternating tensor associated with $c_1(L)$, whose top exterior power is nonzero since $L$ is ample. Intrinsically it is a Hodge map $M:U^\vee\to U(1)$; in the matrix calculation we use the underlying rational spaces. With the Koszul convention for K\"unneth products,
\[
 \operatorname{cl}(E)^2
 =-2\sum_{i<j}(u_i\wedge u_j)\otimes(v_i\wedge v_j).
\]
Poincar\'e duality identifies $H^{2b-2}(D)$ with $(\bigwedge^2U)^\vee$ on underlying rational spaces, by integration against degree-two classes. Consequently $(E^2)_*$ is $-2\bigwedge^2M$ after this identification, and is an isomorphism
\begin{equation}\label{eq:degree-adjustment}
 (E^2)_*:H^{2b-2}(D)(b-2)\xrightarrow{\ \sim\ }H^2(D).
\end{equation}
The twists can also be checked from
$H^{2b-2}(D)(b-2)\simeq H^2(D)^\vee(-2)$ and
$\bigwedge^2M:H^2(D)^\vee\to H^2(D)(2)$.
Composing with $(f,a)_*[Y]$ yields a cycle of codimension $b+2-b=2$ on $S\times D$, nonzero on $V$ by \eqref{eq:albanese-nonzero}. This uses a square of an explicit divisor; it assumes no algebraic inverse Lefschetz operator for a general variety.

Conversely suppose that $\alpha_*:V\to H^2(B)$ is nonzero and $\operatorname{Hg}(V)=\operatorname{SO}(V,q)$. The standard representation is irreducible, so $\alpha_*|_V$ is injective. For $0\ne\sigma\in H^{2,0}(S)$, put $\eta=\alpha_*\sigma\ne0$ in $H^{2,0}(B)$. Write $c=\dim B\ge2$ and let $\ell$ be an ample divisor class on $B$, replaced by a very ample multiple if necessary. The rational surface cycle
\[
 C=\alpha\cdot p_B^*\ell^{c-2}\in\operatorname{CH}_2(S\times B)_{\mathbb Q}
\]
satisfies
\begin{equation}\label{eq:surface-nonzero}
 \int_{S\times B}\operatorname{cl}(C)\cup p_S^*\overline\sigma\cup p_B^*\eta
 =\int_B\overline\eta\cup\eta\cup\ell^{c-2}\ne0.
\end{equation}
Here rationality gives $\alpha_*\overline\sigma=\overline\eta$, and the last pairing is definite up to a fixed nonzero sign on $H^{2,0}(B)$ by the Hodge--Riemann bilinear relations. An actual representative of $C$ is a finite rational combination of integral surfaces; for example use general hyperplanes to represent the indicated intersections. At least one component has nonzero pairing in \eqref{eq:surface-nonzero}. Resolve that component projectively and call the resulting smooth connected surface $Y$. Its projections give morphisms $f:Y\to S$ and $a':Y\to B$ with
\[
 \int_Y f^*\overline\sigma\cup a'^*\eta\ne0.
\]
Thus $f^*\sigma\ne0$, which forces $f$ to have two-dimensional image. It is therefore dominant and generically finite. The projection formula also gives $f_*a'^*\eta\ne0$. Since $\eta$ is a sum of wedges of holomorphic one-forms on $B$, its pullback lies in the wedge image in \eqref{eq:wedge-condition}. This proves the converse.
\end{proof}

\begin{remark}
The reverse implication uses injectivity on $V$, not merely nonvanishing on some divisor class. At a Picard jump a correspondence nonzero only on new algebraic classes need not satisfy the surface condition. No such assertion is used below.
\end{remark}
